\section{서론}
\label{sec:introduction}

대형 언어 모델(LLM)의 파라미터 수가 빠르게 늘면서, 모델 전체를 미세조정하는 길은 점점 멀어졌다. 그 자리를 소수 파라미터만 추가해 모델을 적응시키는 parameter-efficient 방법이 사실상 표준으로 채웠다 \citep{hu2022lora, houlsby2019adapter}. 그중 \textit{soft prompt}는 어휘 바깥의 학습 가능한 연속 임베딩을 입력에 끼워 넣는 도구다. 이 임베딩이 attention을 거치며 모델 행동을 조정한다 \citep{li2021prefix, lester2021power}. 같은 패러다임은 수학 추론을 강화하는 최근 연구에도 폭넓게 쓰인다 \citep{kang2025model, xu2025softcot, ye2025thinking, hao2024training, zhang2025memgen}. 이들 연구는 어휘 밖 연속 임베딩을 \emph{최적화}해 정확도를 끌어올린다는 공통 구조를 공유한다.

그런데 이 효과의 출처가 정말 최적화뿐이라고 단정하기는 어렵다. 학습되지 않은 무작위 연속 임베딩을 그대로 끼워 넣어도 출력 분포는 비자명하게 흔들린다. 한 예로 채팅 형식으로 학습되지 않은 base model에 chat template을 적용하면 형식 불일치 탓에 추론 정확도가 크게 떨어진다. 그런데 무작위 임베딩을 덧붙이면 그 손실이 도리어 회복된다. 단순 노이즈가 모델을 흔드는 데 그친다면 정확도는 더 나빠져야 한다. 따라서 효과의 뿌리는 \emph{학습된 의미}가 아니라 \emph{주입} 자체의 어떤 성질이라는 가설이 자연스럽게 떠오른다. 기존 평가는 최종 정확도에 초점을 맞춰 ``주입''의 효과와 ``최적화''의 효과를 분리하지 못했다. 그래서 어휘 바깥 연속 임베딩이 모델 행동을 어떻게 바꾸는가는 별도의 미해결 문제로 남았다.

본 논문은 학습 없이 주입되는 Random Soft Prompt(RSP)를 분석한다. RSP는 사전학습 임베딩 행렬의 항별 평균과 분산만 따 등방 가우시안에서 추첨한 연속 벡터의 시퀀스다. 무작위성의 역할은 임베딩 분포를 흉내 내는 것이 아니라, 같은 prompt가 rollout마다 \emph{독립}된 섭동을 받도록 보장하는 데 있다. 핵심은 시도 수 자체가 아니라, 매 시도가 초반 토큰의 분기를 다르게 연다는 점이다. 본 논문은 두 질문에 답한다. 첫째, 한 번의 RSP 주입은 왜 \emph{초반} 몇 토큰만 흔들고 나중에는 영향이 사라지는가. 둘째, 시도가 바뀌어 RSP도 새로 뽑힐 때 왜 서로 다른 reasoning path가 시도 사이에서 누적되는가. 첫 질문은 transformer 어텐션 분해가 자연히 만들어 내는 \emph{토큰 단위} 자기조절에서 답을 찾는다. 두 번째 질문은 \emph{N회 독립 재시도}가 작은 분기 확률을 빠르게 누적시키는 사슬에서 답을 찾는다. 한 번의 RSP는 그 시도 안에서 고정된 perturbation으로 작동해 baseline과는 다른 한 갈래를 결정한다. 시도가 누적되면 baseline에서 거의 보이지 않던 후보가 일부 시도에서 답 후보권 안으로 들어오는 사건이 분포로 쌓인다. 마지막으로 이 다양성이 GRPO \citep{shao2024deepseekmath} 학습에 어떤 함의를 갖는지 짚는다.

본 논문의 기여는 다음과 같다.
\begin{itemize}
  \item RSP가 LLM 생성에 어떤 변화를 주는지 실증적으로 분석한다. RSP는 생성 초반의 토큰 엔트로피를 끌어올리고 후반에는 baseline 수준으로 다시 가라앉는다. temperature sampling과 결합하면 더 다양한 추론 궤적이 나온다.
  \item Soft prompt가 학습 없이도 작동하는 이유를 transformer 수준에서 자연어 직관과 함께 정량적으로 설명한다. RSP의 모든 영향은 어텐션 단의 한 가지 양이 통제하고, 그 양은 캐시가 자라면서 자동으로 약해진다. 등방 분포는 어떤 의미 방향에도 사전 가중치를 두지 않으면서, 동시에 모든 방향에 양의 확률을 부여한다.
  \item RSP가 유도하는 다양성이 GRPO 학습에 어떤 함의를 갖는지 논의한다. RSP는 토큰 순위를 \emph{바꾸는} 섭동이고 temperature는 순위를 \emph{보존}하는 섭동이다. 이 차이 덕분에 두 방법은 직교한 탐색 이득을 제공한다.
\end{itemize}
